\cotefile{29}
\scanpage{21}
% archive sheet number: 20
% Formula spaces below reproduce unfilled spaces in the typescript; no missing signs have been supplied.
\deleted{$(a_i)_{i\quad I}$ de $\underline{O}_S$, \unsure{un|une}} une famille $(D_i)_{i\quad I}$ de diviseurs (de Cartier) sur $S$. \unsure{Bour|Pour} tout $S'$ étale sur $S$, on désigne par $\underline{C}_{S'}$ la catégorie suivante: les objets sont les \deleted{\unsure{couples triples}} \inserted{quadruples $M=$} $(\underline{a},\underline{n},\inserted{G,}\underline{u})$, où $\underline{a}=(a_i)_{i\quad I}$ est une famille d'équations pour les $D_{iS'}$, $\underline{n}=(n_i)_{i\quad I}$ une famille d'entiers $1$ premiers aux caractéristiques résiduelles de $S'$, \deleted{\unsure{et} $u:\unsure{G}\quad{}_{\underline{n}}$} $G$ un groupe fini, et $u:G_{S'}\quad{}_{\underline{n}}S'$ un isomorphisme de schémas en groupes sur $S'$. 2tant donné un deuxième tel quadruplet \inserted{$M'=$} $(\underline{a}',\underline{n}',G',u')$, \deleted{un morphisme du second dans le premier est $M'$ dans $M$} $\operatorname{Hom}_{S'}(M',M)$ est vide sauf si $\underline{n}'\quad\underline{n}$, auquel cas c'est l'ensemble des $\underline{b}=(b_i)_{i\quad I}$ tels que $\underline{b}^{n}=\underline{a}'\underline{a}^{-1}$. La composition des morphismes,et la définition des foncteurs changements de base pour une flèche $S''\quad S'$ dans $\underline{\mathrm{Et}}_S$, sont clairs. Le foncteur cartésien $\underline{C}\quad(\mathrm{Sch})/S$ est défini en associant à $M=(\underline{a},\underline{n},G,\underline{u})$ le revêtement Kummérien élémentaire $Z_{\underline{a}}^{\underline{n}}$, \deleted{et à une} \deleted{\unsure{flê|flè}} considéré comme un revêtement à groupe d'opérateurs $G$, grâce à l'isomorphisme $G\quad{}_{\underline{n}}$. On définit de faç\deleted{\unsure{x}}on évidente la valeur de $\underline{r}$ sur les flèches de $\underline{C}_{S'}$, et sa compatibilité avec les changements de base. On voit alors qu'on est sous les conditions de l'exemple 3. De plus, le foncteur $\underline{r}$ est fidèle, et il est pleinement fidèle lorsque $S$ est \deleted{noethérien aux points de} localement noethérien, et de plus normal en les points des $\operatorname{supp}D_i$, et si de plus les $D_i$ sont réduits, et toute \inserted{composante irréductible \unsure{de'un|d'un} $D_i$ est distincte de toute comp. d'un $D_{\unsure{i|j}}$ pour $j\ne i$}

b) Supposons donné un diviseur $D$ sur $S$, \deleted{une famille couvrante $\unsure{S'|S}\quad S$} et considérons le faisceau \unsure{éur|śur|sur} $\underline{\mathrm{Et}}_S$ associé au préfaisceau qui, à tout $S'$, associe l'ensemble des \deleted{diviseurs} ensembles \deleted{$\unsure{\underline{Q}|\underline{\Omega}}$} \inserted{$(D_i)$} $\unsure{\theta|0}$ de diviseurs sur $S'$ ayant les propriétés suivantes: $D_{\unsure{S'|S''}}=\quad D_i$, et pour tout $s'\quad S'$, les restrictions \deleted{non nulles} à $\text{Sêc}(\underline{O}_{S',s'})$ des $D_i$ qui sont non nulles, sont deux à deux disjointes. Soit donnée une section du
